\section{Векторное пространство. Линейные комбинации. Базис и размерность. Модули. Примеры.}

\subsection*{Векторное пространство}

\paragraph{Определение}
Пусть $\mathcal{F} = (F; +, \cdot)$ — поле (элементы — скаляры), а $\mathcal{V} = (V; +)$ — абелева группа (элементы — векторы). $\mathcal{V}$ называется \textbf{векторным пространством} над полем $\mathcal{F}$, если определена операция умножения вектора на скаляр $F \times V \to V$, удовлетворяющая аксиомам ($\forall a, b \in F, \forall x, y \in V$):
\begin{enumerate}
    \item $(a + b)x = ax + bx$
    \item $a(x + y) = ax + ay$
    \item $(a \cdot b)x = a(bx)$
    \item $1x = x$
\end{enumerate}

\paragraph{Теорема (Свойства нулей)}
1. $0x = \mathbf{0}$ (умножение нулевого скаляра на вектор дает нуль-вектор).
2. $a\mathbf{0} = \mathbf{0}$ (умножение скаляра на нуль-вектор дает нуль-вектор).

\textit{Доказательство:}
1. $0x = (1 - 1)x = 1x - 1x = x - x = \mathbf{0}$.
2. $a\mathbf{0} = a(x - x) = ax - ax = \mathbf{0}$. $\square$

\subsection*{Линейные комбинации и независимость}

\paragraph{Определения}
\begin{itemize}
    \item \textbf{Линейная комбинация} векторов $x_i$: сумма вида $\sum_{i=1}^n a_i x_i$, где $a_i \in F$.
    \item \textbf{Линейная независимость:} Множество векторов $\{x_1, \dots, x_k\}$ называется линейно независимым, если равенство $\sum a_i x_i = \mathbf{0}$ выполняется только при всех $a_i = 0$. В противном случае множество \textbf{линейно зависимо}.
\end{itemize}

\paragraph{Теорема}
Линейно независимое множество векторов не содержит нуль-вектора.
\textit{Доказательство:} Если $x_1 = \mathbf{0}$, то можно взять $a_1 = 1$, а остальные $a_i = 0$. Тогда $1 \cdot \mathbf{0} + 0 \cdot x_2 + \dots = \mathbf{0}$, при этом $a_1 \neq 0$, что означает линейную зависимость. $\square$

\subsection*{Базис и размерность}

\paragraph{Определение}
\begin{itemize}
    \item \textbf{Порождающее множество (система образующих):} Подмножество $S$, такое что любой вектор пространства можно представить как линейную комбинацию элементов из $S$.
    \item \textbf{Базис:} Линейно независимое порождающее множество.
    \item \textbf{Размерность:} Количество элементов в базисе (обозначается $\dim V$).
\end{itemize}

\paragraph{Теорема 1 (Единственность разложения)}
Каждый вектор имеет единственное представление в данном базисе.
\textit{Доказательство:} Если $x = \sum a_i E_i$ и $x = \sum b_i E_i$, то $\sum (a_i - b_i) E_i = \mathbf{0}$. В силу линейной независимости базиса $a_i - b_i = 0$, т.е. $a_i = b_i$. $\square$

\paragraph{Теорема 2 (О размере)}
Размер любого линейно независимого множества не превосходит размера базиса ($m \le n$).

\subsection*{Модули}
Если в определении векторного пространства заменить поле $\mathcal{F}$ на произвольное \textbf{кольцо} (не обязательно поле), то полученная структура называется \textbf{модулем} над кольцом. Главное отличие: в модуле не всегда существует базис и не для каждого элемента есть обратный скаляр.

\subsection*{Примеры}
\begin{enumerate}
    \item $F^n$ — пространство кортежей длины $n$ над полем $F$.
    \item \textbf{Булеан:} $(2^M, \Delta)$ является векторным пространством над полем $\mathbb{Z}_2$ (где $1X = X, 0X = \varnothing$).
    \item $V_X$ — пространство формальных сумм $\sum \lambda_x x$, «натянутое» на конечное множество $X$.
    \item Пространство функций $\Phi: X \to F$ с покомпонентными операциями.
\end{enumerate}